% from ega4-16-fr.tex
\oldpage[IV-4]{51}
évidemment de prouver que la suite des fibres des faisceaux figurant dans
(16.9.13.1), en un point $x\in Y'$, est exacte. Or, si l'on pose
$A=\mathcal{O}_{X,x}$, on peut écrire
$\mathcal{O}_{Y,x}=A'=A/\mathfrak{I}$ où $\mathfrak{I}$ est un idéal contenu
dans l'idéal maximal de $A$, puis
$\mathcal{O}_{Y',x}=A'/\mathfrak{K}'$, où $\mathfrak{K}'$ est engendré par une
suite $A'$-régulière d'éléments de $A'$, eux-mêmes images des éléments d'une
suite $A'$-régulière d'éléments de $A$ appartenant à l'idéal maximal de $A$.
Si $\mathfrak{K}$ est l'idéal engendré par ces derniers et
$\mathfrak{L}=\mathfrak{I}+\mathfrak{K}$, on a
$\mathcal{O}_{Y',x}=A/\mathfrak{L}$, et comme on est dans la situation de
(16.9.13.5), l'homomorphisme canonique
\[
  \mathfrak{I}/\mathfrak{I}\mathfrak{L}
  \longrightarrow(\mathfrak{I}+\mathfrak{L}^2)/\mathfrak{L}^2
\]
est bijectif. Mais cela montre que la suite
\[
  0\longrightarrow\mathfrak{I}/\mathfrak{I}\mathfrak{L}
  \longrightarrow\mathfrak{L}/\mathfrak{L}^2
  \longrightarrow(\mathfrak{L}/\mathfrak{I})/
    (\mathfrak{L}/\mathfrak{I})^2\longrightarrow0
\]
est exacte (voir la démonstration de (16.2.7)), et les modules figurant dans
cette suite sont précisément les fibres en $x$ des faisceaux de (16.9.13.1).
La seconde assertion résulte de ce que $\mathcal{N}_{Y'/Y}$ est un
$\mathcal{O}_{Y'}$-Module localement \emph{libre} (16.9.8) et de Bourbaki,
\emph{Alg.}, chap.~II, 3\textsuperscript{e} éd., §~1, n\textsuperscript{o}~11,
prop.~21.

\subsection{Morphismes différentiellement lisses}
\EGAReaderTarget{ega4.c04.subsec.16-10}{IV.16.10}

\begin{definition}[16.10.1]
\EGAReaderTarget{ega4.c04.def.16-10-1}{IV.16.10.1}
On dit qu'un morphisme de préschémas $f:X\to S$ est
\emph{différentiellement lisse} (ou que $X$ est \emph{différentiellement
lisse sur $S$}) s'il vérifie les conditions suivantes:
\begin{enumerate}
  \item[(i)] $\Omega_{X/S}^1$ est un $\mathcal{O}_X$-Module localement
  projectif, c'est-à-dire que tout point de $X$ admet un voisinage ouvert
  affine $U$ tel que $\Gamma(U,\Omega_{X/S}^1)$ soit un
  $\Gamma(U,\mathcal{O}_X)$-module projectif (non nécessairement de type
  fini).
  \item[(ii)] L'homomorphisme canonique (16.3.1.1)
  \[
    \mathbf{S}_{\mathcal{O}_X}^{\bullet}(\Omega_{X/S}^1)
    \longrightarrow\mathscr{G}r_{\bullet}(\mathscr{P}_{X/S})
  \]
  est bijectif.
\end{enumerate}
En particulier, si $\Omega_{X/S}^1$ est localement libre de rang fini, les
$\mathscr{P}_{X/S}^n$ sont des $\mathcal{O}_X$-Modules localement libres de
rang fini (étant des extensions de tels Modules).
\end{definition}

On dit que $f$ est \emph{différentiellement lisse en un point $x\in X$} (ou
que $X$ est différentiellement lisse sur $S$ au point $x$) s'il existe un
voisinage ouvert $U$ de $x$ dans $X$ tel que $f|U$ soit différentiellement
lisse.

Nous verrons plus loin (17.12.4) qu'un morphisme lisse est différentiellement
lisse, ce qui justifie la terminologie; mais la réciproque est inexacte; en
effet, un \emph{monomorphisme} $f:X\to S$ est différentiellement lisse,
puisque $\Omega_{X/S}^1=0$ en vertu de (I, 5.3.8), et par suite
l'homomorphisme surjectif (16.3.1.1) est évidemment bijectif; or un
monomorphisme n'est même pas nécessairement plat, ni \emph{a fortiori} lisse.
Bornons-nous ici à noter la proposition suivante:

\begin{proposition}[16.10.2]
\label{IV.16.10.2}
Soient $A$ un anneau, $B$ une $A$-algèbre formellement lisse pour les
topologies discrètes (0, 19.3.1). Alors $\operatorname{Spec}(B)$ est
différentiellement lisse sur $\operatorname{Spec}(A)$.
\end{proposition}

En effet, $B\otimes_A B$ est alors (pour les topologies discrètes) une
$B$-algèbre formellement lisse (pour l'un ou l'autre des homomorphismes
canoniques $b\mapsto b\otimes1$, $b\mapsto1\otimes b$ de $B$
